\chapter{Sprint 3 : Gestion des tâches et Kanban}

\section*{Introduction}
Le troisième sprint constitue le cœur opérationnel de la plateforme. Il introduit
la gestion hiérarchique des tâches (avec sous-tâches), les dépendances de type
FS/SS/FF/SF avec détection de cycles (BFS), l'affectation des membres avec
allocation en pourcentage, le tableau Kanban drag-and-drop avec mise à jour
optimiste, et la saisie du temps passé.

% ============================================================
\section{Sprint Backlog}

\begin{table}[H]
  \centering
  \caption{Sprint Backlog — Sprint 3}
  \label{tab:backlog_s3}
  \begin{tabularx}{\textwidth}{C{1.2cm}L{5cm}L{2cm}C{2cm}C{1.5cm}}
    \toprule
    \thead{ID} & \thead{User Story} & \thead{Acteur} & \thead{Priorité} & \thead{Points} \\
    \midrule
    \rowcolor{lightgray}
    US-11 & Créer des tâches hiérarchiques (sous-tâches)     & Member & Must   & 8 \\
    US-12 & Définir des dépendances FS/SS/FF/SF              & Member & Must   & 8 \\
    \rowcolor{lightgray}
    US-13 & Affecter des membres avec allocation \%          & Member & Must   & 5 \\
    US-14 & Déplacer une tâche en Kanban (drag-and-drop)     & Member & Must   & 5 \\
    \rowcolor{lightgray}
    US-15 & Saisir le temps passé sur une tâche             & Member & Should & 3 \\
    US-16 & Commenter une tâche et mentionner (@)            & Member & Should & 5 \\
    \midrule
    \multicolumn{4}{r}{\textbf{Total :}} & \textbf{34} \\
    \bottomrule
  \end{tabularx}
\end{table}

% ============================================================
\section{Conception}

\subsection{Diagramme de cas d'utilisation}

\begin{figure}[H]
  \centering
  \IfFileExists{figures/sprint3_use_case.png}{%
    \includegraphics[width=0.88\textwidth]{figures/sprint3_use_case.png}%
  }{%
    \fbox{\parbox{0.88\textwidth}{\centering\vspace{1.8cm}
      \textbf{[IMAGE A INSÉRER]}\\[0.4em]
      \texttt{figures/sprint3\_use\_case.png}\\[0.3em]
      \textit{Cas d'utilisation : Member, Chef. Dépendances, détection cycle, export CSV}
    \vspace{1.8cm}}}%
  }
  \caption{Diagramme de cas d'utilisation — Sprint 3}
  \label{fig:uc_s3}
\end{figure}

\subsection{Diagramme de séquence — Déplacement Kanban}

\begin{figure}[H]
  \centering
  \IfFileExists{figures/sprint3_sequence_kanban.png}{%
    \includegraphics[width=0.90\textwidth]{figures/sprint3_sequence_kanban.png}%
  }{%
    \fbox{\parbox{0.90\textwidth}{\centering\vspace{1.8cm}
      \textbf{[IMAGE A INSÉRER]}\\[0.4em]
      \texttt{figures/sprint3\_sequence\_kanban.png}\\[0.3em]
      \textit{Séquence drag-and-drop Kanban : mise à jour optimiste → PATCH API → broadcast Reverb}
    \vspace{1.8cm}}}%
  }
  \caption{Diagramme de séquence — Déplacement d'une tâche sur le Kanban}
  \label{fig:seq_kanban}
\end{figure}

Le déplacement d'une tâche se déroule en six étapes : (1) drag-and-drop de la
carte, (2) mise à jour optimiste de l'interface React, (3) \texttt{PATCH
/api/tasks/\{id\}/move}, (4) vérification par \texttt{CheckProjectAccess}, (5)
mise à jour en base et diffusion via Laravel Reverb, (6) les autres clients
connectés reçoivent l'événement WebSocket et mettent à jour leur vue. En cas
d'erreur serveur, la mise à jour optimiste est annulée (\textit{rollback} React
Query).

\subsection{Diagramme de classes}

\begin{figure}[H]
  \centering
  \IfFileExists{figures/sprint3_class.png}{%
    \includegraphics[width=0.88\textwidth]{figures/sprint3_class.png}%
  }{%
    \fbox{\parbox{0.88\textwidth}{\centering\vspace{1.8cm}
      \textbf{[IMAGE A INSÉRER]}\\[0.4em]
      \texttt{figures/sprint3\_class.png}\\[0.3em]
      \textit{Classes : Task, TaskAssignment, TaskDependency, TaskComment, Attachment, TimeEntry}
    \vspace{1.8cm}}}%
  }
  \caption{Diagramme de classes — Sprint 3}
  \label{fig:class_s3}
\end{figure}

% ============================================================
\section{Réalisation}

\subsection{Vue Kanban}

\begin{figure}[H]
  \centering
  \IfFileExists{figures/sprint3_14_kanban.png}{%
    \includegraphics[width=0.90\textwidth]{figures/sprint3_14_kanban.png}%
  }{%
    \fbox{\parbox{0.90\textwidth}{\centering\vspace{2cm}
      \textbf{[IMAGE A INSÉRER]}\\[0.4em]
      \texttt{figures/sprint3\_14\_kanban.png}\\[0.3em]
      \textit{Tableau Kanban : colonnes À faire / En cours / À risque / Terminé, drag-and-drop}
    \vspace{2cm}}}%
  }
  \caption{Tableau Kanban — Drag-and-drop avec mise à jour temps réel}
  \label{fig:ui_kanban}
\end{figure}

\subsection{Formulaire wizard 4 étapes}

\begin{figure}[H]
  \centering
  \begin{subfigure}[b]{0.48\textwidth}
    \IfFileExists{figures/sprint3_15_task_step1.png}{%
      \includegraphics[width=\textwidth]{figures/sprint3_15_task_step1.png}%
    }{%
      \fbox{\parbox{\textwidth}{\centering\vspace{1.2cm}\small
        \texttt{sprint3\_15\_task\_step1.png}\\Étape 1 : Infos de base\vspace{1.2cm}}}%
    }
    \caption{Étape 1 : Infos de base}
  \end{subfigure}
  \hfill
  \begin{subfigure}[b]{0.48\textwidth}
    \IfFileExists{figures/sprint3_16_task_step2.png}{%
      \includegraphics[width=\textwidth]{figures/sprint3_16_task_step2.png}%
    }{%
      \fbox{\parbox{\textwidth}{\centering\vspace{1.2cm}\small
        \texttt{sprint3\_16\_task\_step2.png}\\Étape 2 : Planification\vspace{1.2cm}}}%
    }
    \caption{Étape 2 : Planification et dates}
  \end{subfigure}
  \caption{Wizard de création de tâche — Étapes 1 et 2}
  \label{fig:ui_wizard_1_2}
\end{figure}

\begin{figure}[H]
  \centering
  \begin{subfigure}[b]{0.48\textwidth}
    \IfFileExists{figures/sprint3_17_task_step3.png}{%
      \includegraphics[width=\textwidth]{figures/sprint3_17_task_step3.png}%
    }{%
      \fbox{\parbox{\textwidth}{\centering\vspace{1.2cm}\small
        \texttt{sprint3\_17\_task\_step3.png}\\Étape 3 : Dépendances\vspace{1.2cm}}}%
    }
    \caption{Étape 3 : Compétences et dépendances}
  \end{subfigure}
  \hfill
  \begin{subfigure}[b]{0.48\textwidth}
    \IfFileExists{figures/sprint3_18_task_step4_ml.png}{%
      \includegraphics[width=\textwidth]{figures/sprint3_18_task_step4_ml.png}%
    }{%
      \fbox{\parbox{\textwidth}{\centering\vspace{1.2cm}\small
        \texttt{sprint3\_18\_task\_step4\_ml.png}\\Étape 4 : Assignation ML\vspace{1.2cm}}}%
    }
    \caption{Étape 4 : Recommandation IA}
  \end{subfigure}
  \caption{Wizard de création de tâche — Étapes 3 et 4}
  \label{fig:ui_wizard_3_4}
\end{figure}

\subsection{Code clé — Détection de cycles (BFS)}

\begin{lstlisting}[style=phpstyle,caption={TaskDependencyService — Détection de cycle BFS},label={lst:cycle_bfs}]
// app/Services/TaskDependencyService.php
public function add(Task $task, int $dependsOn,
                    string $type = 'FS'): void
{
    if ($task->id === $dependsOn) {
        abort(422, 'Cycle direct détecté.');
    }
    if ($this->wouldCreateCycle($task->id, $dependsOn)) {
        abort(422, 'Cycle de dépendances détecté.');
    }
    TaskDependency::create([
        'task_id'             => $task->id,
        'depends_on_task_id'  => $dependsOn,
        'type'                => $type,
    ]);
}

private function wouldCreateCycle(int $taskId, int $dependsOnId): bool
{
    $visited = [];
    $queue   = [$dependsOnId];

    while ($queue) {
        $current = array_shift($queue);
        if ($current === $taskId) return true;
        if (isset($visited[$current])) continue;
        $visited[$current] = true;

        TaskDependency::where('task_id', $current)
            ->pluck('depends_on_task_id')
            ->each(fn($d) => $queue[] = $d);
    }
    return false;
}
\end{lstlisting}

% ============================================================
\section{Tests unitaires}

\begin{lstlisting}[style=phpstyle,caption={Tests unitaires — Dépendances et Kanban},label={lst:tests_s3}]
it('detects direct cycle', function () {
    $task = Task::factory()->create();
    expect(fn() => app(TaskDependencyService::class)->add($task, $task->id))
        ->toThrow(HttpException::class, 422);
});

it('detects indirect cycle A -> B -> C -> A', function () {
    [$a, $b, $c] = Task::factory()->count(3)->create();
    app(TaskDependencyService::class)->add($b, $a->id);
    app(TaskDependencyService::class)->add($c, $b->id);
    expect(fn() => app(TaskDependencyService::class)->add($a, $c->id))
        ->toThrow(HttpException::class, 422);
});

it('moves task to new kanban column', function () {
    $task   = Task::factory()->create(['workflow_status' => 'todo']);
    $status = TaskStatus::factory()->create(['project_id' => $task->project_id]);
    $this->actingAs($task->project->owner)
         ->patchJson("/api/tasks/{$task->id}/move", [
             'status_id'       => $status->id,
             'workflow_status' => 'in_progress',
         ])->assertOk()
           ->assertJsonFragment(['workflow_status' => 'in_progress']);
});
\end{lstlisting}

\section*{Conclusion}
Le wizard 4 étapes a significativement amélioré la qualité de saisie : les champs
obligatoires sont groupés en étape~1, les dates avec validations en étape~2, les
compétences et dépendances en étape~3, et la recommandation IA en étape~4. La
détection de cycles BFS protège l'intégrité des données. Le Kanban drag-and-drop
avec mise à jour optimiste est, selon les retours utilisateurs, la fonctionnalité
la plus appréciée du sprint.
